\section{The simulator: implementation, information safety, verification, and the v0.6 repairs}
\label{sec:engine}

The simulator is the instrument every number in this report is read from, so its
properties are stated before its results. Three of them matter: it cannot leak a
hidden field to a policy, it checks its own deductions against the truth while it
runs, and a configuration in which every v0.6 mechanism is switched off is
\vfive{} bit for bit. The last of those is not a courtesy. It is what makes the
ablation table of \S\ref{sec:results-ablations} exact rather than approximate,
and it is the first thing \filepath{engine/experiments_v06.sh} runs.

This section also records four repairs to the engine itself. They are reported
here rather than folded silently into a diff because three of them are the
reason a measurement in this report was possible at all, and because one of them
--- a dead command-line option --- means a published pass in the v0.4 record was
vacuous (\S\ref{sec:corr-vfour}).

\subsection{Implementation}
\label{sec:engine-impl}

The engine is unchanged in structure from the v0.4 report, which describes it in
full \cite{fishbot04}: a single-header C++20 core in \filepath{engine/src/}, a
hand as one \verb|uint64_t| over the 54 card indices $c = 6h + i$, and
counter-based random streams keyed by \emph{(seed, seat, purpose)} so that a deal
and its entire play reproduce from the seed alone. The duplicate-block protocol
of \S\ref{sec:protocol} depends on that last property and on nothing else.

What moved in v0.6 is the policy layer. \texttt{V06Agent} derives from
\texttt{V05Agent} and defers to it whenever no v0.6 mechanism is enabled, so the
v0.5 decision path is not reimplemented but inherited; the determinized search of
\S\ref{sec:search} and the rollout blueprint it calls live in two new headers,
\filepath{engine/src/v06.hpp} and \filepath{engine/src/v06_rollout.hpp}. The
diagnostics this report leans on are a new subcommand rather than scratch
probes: \texttt{fish v6probe --mode=ties} censuses ask decisions and every
tie-break rule including hindsight, \texttt{--mode=belief} scores several
posteriors as predictors on identical states, and \texttt{--mode=search} reports
how often a guarded search actually deviates from its blueprint. All three are in
\filepath{engine/src/probe_v06.hpp} and all three are run by the battery, so the
figures they produce are regenerated rather than transcribed.

Throughput on the shipped configuration is \vsixThroughputSix{} games/s against
\vsixThroughputFive{} for \vfive{} and \vsixThroughputFour{} for \vfour{}, all on
the same machine at \vsixCores{} threads. Every compute figure in this report is
single-machine and was measured under load, so all of them are lower bounds.

\subsection{Information safety}
\label{sec:engine-safety}

The acting policy never receives the game state. It receives its own
\texttt{Knowledge} object --- the constraint bookkeeping of \S\ref{sec:belief}
--- and the public record, in which each event carries only the actor, the
target, the named card, the outcome, any declared allocation, and the resulting
public hand counts. There is no interface through which a policy can read
another player's cards, so the class of leak the \vthree{} audit reported in its
predecessor \cite{fishbot03} is excluded structurally rather than by inspection.
The property and the argument for it are unchanged from the v0.4 report
\cite{fishbot04}.

Soundness is the stronger claim, and it is checked rather than argued. After
every public event the engine's audit mode verifies, for every player and against
the true deal, that every card the constraint state records as resolved is held
by the seat it is resolved to; that the true holder of every unresolved card lies
in that card's surviving support; that the derived capacities equal the true
counts of unresolved cards per player; and that every live ask-legality
certificate is satisfied by the true deal. A single violation would mean the
belief engine had excluded the truth.

\subsection{Verification suite}
\label{sec:engine-verify}

\texttt{fish verify} plays the ordered cross-product of a \numVerifyPool{}-policy
pool --- \numVerifyPairs{} ordered pairs, so both orderings of every pair and
every self-play pairing are exercised --- with the audit above enabled
throughout, and additionally checks deck integrity and card conservation, that
exactly nine half-suits are resolved in every completed game, that a declaration
only ever assigns cards to the declarer's own teammates, that ask legality is
enforced independently on the move generator and on the applier, and that
replaying a seed at a fixed thread count reproduces the game bit-identically. It
returns a non-zero exit status on failure, which is what lets the battery use it
as a gate rather than as a report.

Over \vsixVerifyGames{} deals the suite recorded \vsixAuditViolations{}
violations in \vsixAuditChecks{} checks, with an overall verdict of
\vsixVerify{}. The random streams are keyed so that thread count cannot affect a
game's outcome, but the suite tests replay identity only at a fixed thread count
and does not test that claim directly; this is carried forward from the v0.4
study as an unclosed gap and is listed as one in \S\ref{app:repro-gaps}.

\subsection{The v0.6 engine repairs}
\label{sec:engine-repairs}

Each repair below is stated as a defect, the measurement that shows the defect is
real, what the defect blocked, and the measurement that shows it is fixed. Three
of the four are correctness-neutral for \vfive{}: they change nothing a \vfive{}
game does, which is why they survived a release. The fourth changes the Fast
posterior and is therefore a flag that defaults off.

\paragraph{BlockDP tables aliased across instances.}
\filepath{engine/src/blockdp.hpp} parked every instance's Unit, Group, forward
and backward tables in one \verb|thread_local| buffer pool, and
\texttt{BlockDP::build} took its storage from that pool without owning it. Two
\texttt{BlockDP} objects alive on one thread therefore shared storage, and a
second instance's \texttt{build} silently repointed the first instance's tables
at its own data. The driver makes this routine rather than exotic: agents own one
\texttt{block} member each and rebuild only when dirty, while
\texttt{Game::forcedEndgame} and \texttt{Game::declarationRound} poll several
agents back to back between public events.

The defect is real and was measured directly. Reading one instance's own group
table twice with an unrelated \texttt{build} in between gave
\numAliasMismatch{} mismatches in \numAliasChecks{} checks: one observer's
posterior tables had been replaced by another observer's
(\filepath{research/v05/results/P2-forced-endgame.md}, \S6). An earlier form of
the same check that issued a real query after the clobber died with
\texttt{SIGBUS}, because the stale forward and backward pointers index a table
sized for a different capacity vector.

What it blocked is the whole of \S\ref{sec:belief}. The deployed approximate
belief never builds a \texttt{BlockDP}, so nothing in the v0.4 or v0.5 release
path could observe the defect; anything exact hits it on the first
multi-agent poll. That asymmetry is exactly how a fatal defect survives into a
release, and it is why the exact posterior could not be evaluated as a predictor
before this study.

The repair is a per-thread generation stamp plus a stored copy of the source
\texttt{Knowledge}. Every \texttt{build} increments the counter and records its
value; every public query calls \texttt{ensureCurrent}, which re-derives the
instance lazily from its stored \texttt{Knowledge} when it finds its own stamp
stale. The check cannot be forgotten at a call site because it lives inside the
queries, and it costs $O(1)$ memory. The probe in
\filepath{engine/src/probe_forcedendgame.hpp} was extended to separate the two
things the original check conflated: \texttt{QUERY} mismatches, where the same
query issued on the same instance before and after an unrelated build returns
different answers, and raw shared-pool field reads, where the instance's own
pointers still address the pool and the fields behind them have changed. The
first must be zero and is \vsixAliasQueryMismatch{}. The second is
\vsixAliasRawMismatch{} and is \emph{not} fixed: the pool is still shared, the
raw fields still change, and the guard's contract is that no public query ever
reads them stale. Reporting only the first would overstate the repair.

\paragraph{\texttt{gateaudit} was dead code for v0.5.}
The declaration pre-gate audit is instrumentation that counts how many
declaration opportunities a cheap pre-gate rejects and how many of those
rejections were false negatives --- candidates the expensive test would have
accepted. The switch that enables it was parsed only in the \vfour{} branch of
\filepath{engine/src/factory.hpp}. A v0.5 specification therefore parsed
\texttt{gateaudit=1}, ignored it, and reported a pass over zero opportunities:
\texttt{fish gateaudit --a=v05:gateaudit=1} could not fail, because it never
looked at anything. \S\ref{sec:corr-vfour} records what that means for the v0.4
record's corresponding claim.

The option is now parsed in the v0.5 and v0.6 branches as well. The audit over
\vsixGateOpportunities{} declaration opportunities rejected \vsixGateRejected{}
at the cheap gate with \vsixGateFalseNeg{} false negatives, verdict
\vsixGateVerdict{}. The instrument is only as good as its opportunity count, so
that count is reported first and the verdict second.

\paragraph{Two defects in the tuner.}
\filepath{engine/src/tuner.hpp} chose between the two parameter-passing spellings
with \verb|w.size() > 18| where the ask-feature count \verb|NFEAT| ---
\numAskFeats{} --- belongs. The expression is correct today by coincidence, and
it is the same class of defect, in the same place, that cost the v0.4 study a
fitting round when the ask-feature count grew. It now reads \verb|NFEAT|. The
second defect is worse because it is silent and quantitative: the win rate was
computed as \verb|st.winsA / (st.games * 2)|, hard-coding the tuner's two
rotations into the denominator. Fitting at six rotations --- which the
duplicate-block design of \S\ref{sec:protocol} makes the natural choice --- would
have reported a third of the true win rate for every candidate, with no error and
no warning. It now reads \verb|st.games * mc.rotations|.

\paragraph{Certificate de-duplication, opt-in.}
\texttt{Knowledge::onEvent} in \filepath{engine/src/belief.hpp} appended a fresh
ask-legality certificate on every repeated ask without checking whether an
identical one, or one strictly stronger, was already stored. The store therefore
grows carrying no information, and every consumer that iterates it pays for the
duplicates. The obvious repair --- skip a certificate that duplicates a stored
one, or whose card set is a superset of a stored one for the same player --- is
implemented, and it is a flag that defaults off. The reason is that the
duplicates are not inert. The approximate conditioning step of
\S\ref{sec:belief} applies a certificate once per stored copy, so removing
copies \emph{changes} the Fast posterior and therefore changes play. Making
de-duplication the default would have silently altered every \vfour{} and
\vfive{} number this report compares against. With the flag off, both stay bit
for bit what they were.

\paragraph{The identity control, and why it is a result.}
\texttt{v06:legacy=1} restores \vfive{}'s parameter vector and zeroes the three
v0.6 ask terms, and with every mechanism switch off \texttt{V06Agent} defers to
its base class. The battery checks that this is an identity and not a close
agreement, by comparing the md5 digest of the entire \texttt{fish pathology}
output --- every pathology counter, not a win rate --- between
\texttt{v06:legacy=1} and \texttt{v05} on the same deals and seed. The verdict is
\vsixIdentity{}.

This matters more than a build-hygiene check. Every ablation in
\S\ref{sec:results-ablations} and \S\ref{app:ablations} is a paired comparison
between the v0.6 binary carrying mechanism $X$ and the v0.6 binary carrying
nothing. If the all-switches-off configuration merely resembled \vfive{}, each
row would confound the mechanism with whatever else differed, and the table would
be an approximation whose error had no bound. Because the configuration is
\vfive{} bit for bit, the difference each row reports is the mechanism and
nothing else. The control is cheap, it is run before any strength number is
collected, and it is the reason the ablation table can be read as exact.
